\usepackage{tikz}
\usepackage{pgfplots}
\usepackage{pgfplotstable}
\pgfplotsset{compat=1.8}

%%%%%%%%%%%%%%%%%%%%%%%%%%%%%%%%%%%%%%%%%%%%%%%%%%%%%%%%%%%%%%%%%%%%%%%%%%%%%%%
%
% Annotated demo macros for convergence plots.
%
% Some info:
% - lines starting with # are ignored
% - by default pgfplotstable expects the first line to be a header
% - no separator between column header labels
% - by default pgfplotstable uses whitespace characters as column separator
%
% Further questions: <aurelien.larcher@gmail.com>
%
%%%%%%%%%%%%%%%%%%%%%%%%%%%%%%%%%%%%%%%%%%%%%%%%%%%%%%%%%%%%%%%%%%%%%%%%%%%%%%%

\def\plotspecs{licorne}
\def\timesteplabel{TimeStep}
\def\meshsizelabel{MeshSize}

\def\suffixDim{:dim}
\newcommand{\spaceDim}[1]{#1\suffixDim}

\def\suffixSol{:sol}
\newcommand{\normSol}[2]{#1\suffixSol #2}

\def\suffixErr{:err}
\newcommand{\normErr}[2]{#1\suffixErr #2}

%------------------------------------------------------------------------------
% Default axis style
%------------------------------------------------------------------------------
\pgfplotsset{
ylabel style={overlay},
yticklabel style={overlay},
every axis/.append style={
    label style={font=\scriptsize},
    tick label style={font=\scriptsize},
    legend style={font=\tiny}
}}

% -----------------------------------------------------------------------------
% Encapsulate table loading
% #1: shape index N (number of batches)
% #1: shape index M (number of rows per batch)
% -----------------------------------------------------------------------------

\newcommand{\loadtable}[2] {
    \pgfplotstableread[]{#1}{#2}
}

% -----------------------------------------------------------------------------
% Style to select only points in given range:
% #1: start index
% #2: end index (inclusive)
% -----------------------------------------------------------------------------

\pgfplotsset{
    select coords between index/.style 2 args= {
        x filter/.code= {
            \ifnum\coordindex<#1\def\pgfmathresult{}\fi
            \ifnum\coordindex>#2\def\pgfmathresult{}\fi
        }  
    }
}

% -----------------------------------------------------------------------------
% Style to select only points in given offset and modulus:
% #1: shape index N
% #2: shape index M (modulus)
% -----------------------------------------------------------------------------

\pgfplotsset{
    select coords modulo index/.style 2 args= {
        x filter/.code= {
            % Technically this should work but it fails to compile
            %\pgfmathparse{int(mod(\coordindex,#2))}
            %\pgfmathtruncatemacro{\k}{\pgfmathresult}
            \newcount\k
            \k=\coordindex
            \divide\k by #2
            \multiply\k by #2
            \multiply\k by -1
            \advance\k by \coordindex\relax
            \the\k
            \ifnum\k=#1\relax\else\def\pgfmathresult{}\fi
            }  
    }
}

% -----------------------------------------------------------------------------
% Plot data in range (a, b)
% #1: data handle
% #2: a
% #3: b
% #4: column label for abscissa
% #5: column label for ordinate
% -----------------------------------------------------------------------------

\newcommand{\rangeplot}[7] {
    \addplot plot
        table [x = #4, y = #5, select coords between index={#2}{#3}] {#1};
    \if\relax\detokenize{#7}\relax
			\addlegendentryexpanded{{#6}}\relax
    \else
    	\pgfplotstablegetelem{#2}{#7}\of#1
    	\pgfmathsetmacro{\hn}{\pgfplotsretval}
    	\addlegendentryexpanded{${#6}_\curve=\hn$}\fi
}

% -----------------------------------------------------------------------------
% Plot data by chunk with given (x, y) columns
% #1: data handle
% #2: shape index N
% #3: shape index M
% #4: column label for abscissa
% #5: column label for ordinate
% -----------------------------------------------------------------------------

\newcommand{\multiplot}[7] {
    \newcount\a
    \newcount\b
    \foreach \curve in {1,...,#2} {
      \pgfmathsetcount{\a}{#3 * (\curve - 1)}
      \pgfmathsetcount{\b}{\a + #3 - 1}
      \addplot plot
          table [x = #4, y = #5, select coords between index={\a}{\b}] {#1};
      \pgfplotstablegetelem{\a}{#7}\of#1
      \pgfmathsetmacro{\hn}{\pgfplotsretval}
      \addlegendentryexpanded{${#6}_\curve=\hn$}
    }
}

% -----------------------------------------------------------------------------
% Plot data by set range set with given (x, y) columns
% #1: data handle
% #2: offset
% #3: shape index M
% #4: column label for abscissa
% #5: column label for ordinate
% -----------------------------------------------------------------------------

\newcommand{\moduloplot}[6] {
    \addplot plot
        table [x = #4, y = #5, select coords modulo index={#2}{#3}] {#1};
    \if\relax\detokenize{#6}\relax
			\else\addlegendentryexpanded{{#6}}\relax\fi
    
    %\addplot [orange, select coords modulo index={#2}{#3}]
    %    table [x = #4, y ={create col/linear regression={y= #5}}]{#1};
}

% -----------------------------------------------------------------------------
% Plot time convergence (shortcut to et TimeStep as default)
% #1: data handle
% #2: shape index N
% #3: shape index M
% #4: column label
% -----------------------------------------------------------------------------

\newcommand{\timeconvergence}[4] {
    \multiplot{#1}{#2}{#3}{\timesteplabel}{#4}{h}{\meshsizelabel}
}

% -----------------------------------------------------------------------------
% Example: Plot for $foo, L2 norm ($foo:errL2) and H1 seminorm ($foo:errH1)
% #1: data handle
% #2: shape index N
% #3: shape index M
% #4: function label
% -----------------------------------------------------------------------------

\newcommand{\timeLtwoHone}[4] {
    \begin{minipage}[t]{0.5\textwidth}
        \begin{tikzpicture}
            \begin{loglogaxis}[xlabel=$dt$, ylabel=#4 $\xLtwo$ Error, legend
            pos=north west]
                \timeconvergence{#1}{#2}{#3}{\normErr{#4}{L2}}
            \end{loglogaxis}
        \end{tikzpicture}
    \end{minipage}
    \begin{minipage}[t]{0.5\textwidth}
        \begin{tikzpicture}
            \begin{loglogaxis}[xlabel=$dt$, ylabel=#4 $\xHone_0$ Error, legend
            pos=north west]
                \timeconvergence{#1}{#2}{#3}{\normErr{#4}{H1}}
            \end{loglogaxis}
        \end{tikzpicture}
    \end{minipage}
}

%------------------------------------------------------------------------------
